Grounding metrics test whether the model is using the formal vocabulary of the domain and problem. They are not success metrics by themselves as a model can use legal names and still fail.\\
Their role is discriminatory because they separate models that at least read the scheme from models that invent actions or objects.

\begin{table}[H]\small
\centering
\begin{tabularx}{\linewidth}{@{}l l Y Y@{}}
\toprule
\textbf{Metric} & \textbf{Formula} & \textbf{What it measures} & \textbf{Main interpretation}\\
\midrule
Action hallucination & \(\frac{\mathrm{illegal\ action\ names}}{\mathrm{total\ actions}}\) & Emitted action names not present in the schema. & High value means weak action grounding.\\
\addlinespace
Fuzzy hallucination & Strict hallucinations with edit distance \(>2\). & Genuine inventions versus spelling or formatting errors. & A strict/fuzzy gap suggests parser or formatting issues.\\
\addlinespace
Object hallucination & \(\frac{\mathrm{illegal\ arguments}}{\mathrm{total\ arguments}}\) & Arguments not present in problem objects. & Catches legal verbs applied to invented entities.\\
\addlinespace
IHR & \(1-\mathrm{action\ hallucination}\) & Grounding score with higher values better. & Used in the composite score and profiles.\\
\bottomrule
\end{tabularx}
\caption{Grounding and discriminatory metrics.}
\label{tab:grounding-metrics-compact}
\end{table}

The grounding layer is a gate for later interpretation. If a model hallucinates many action names, execution and CoT metrics become less reliable because there is no stable formal plan to inspect. Low hallucination only means the model is inside the legal vocabulary but it does not imply planning ability.